\section{Discussion}

\subsection{Key Findings and Their Interpretation}

Our results confirm the three-step mechanism at three independent measurement levels. The distributional evidence (h-m1) establishes why the pattern exists; the rank-correlation evidence (h-m2) confirms it without threshold assumptions; the AUROC evidence (h-m3) quantifies the practical magnitude. This layered structure provides stronger mechanistic support than any single metric and is unusual for hallucination detection papers.

The symmetric magnitude of P1 and P2 ($\approx$0.10--0.12 AUROC on LLaMA across both benchmark types) is particularly noteworthy. Both effects are approximately equal in magnitude but opposite in direction --- consistent with a single underlying distribution-shape driver rather than two separate phenomena.

The cross-architecture consistency (LLaMA-2-7B and Mistral-7B-v0.1) within the 7B parameter scale is the strongest available generalizability evidence. These models differ in architecture, pretraining corpus, and tokenization; consistent directional effects across both argue against model-specific confounds.

\subsection{The Raw\_sum Discovery}

The refutation of P3 --- raw\_sum is the strongest, not weakest, aggregation on TriviaQA --- is the most practically impactful result. The most likely explanation is that the ``length bias'' is a signal: on factual recall benchmarks, correct answers tend to be longer and each token is high-confidence. The unnormalized sum accumulates a larger negative value for correct answers, making it more discriminative than per-token statistics. This interpretation is consistent with raw\_sum being weakest on TruthfulQA (where all tokens are uniformly confident regardless of correctness).

We note this requires length-stratified validation. The h-e1 codebase includes a \texttt{length\_stratified\_auroc} function implemented but not yet executed.

\subsection{Comparison to Published Baselines}

Our min log-prob on TriviaQA (0.849--0.892) exceeds published CCP/mean log-prob (0.72--0.80) by 0.07--0.09 AUROC. Our raw\_sum (0.895--0.896) exceeds Semantic Entropy ($\approx$0.79) by $\approx$0.10 AUROC with 10$\times$ lower inference cost. These cross-pipeline comparisons should be interpreted with appropriate caution --- different hardware, dataset splits, tokenization, and evaluation protocols may account for part of the gap --- but the magnitude of the gap makes a purely methodological explanation unlikely.

\subsection{Limitations}

\textbf{NQ data gap.} P1 is confirmed on TriviaQA only; NQ inference ran but cache was not persisted. The mechanism prediction for NQ is strong (same hallucination type), and re-running requires only GPU time.

\textbf{7B model scale.} All results are at 7B parameters. TruthfulQA inverse scaling \citep{lin2022} predicts stronger P2 effects at 70B+; our 7B P2 results may underestimate the mechanism.

\textbf{Greedy decoding only.} All results are valid within the greedy-decoding regime; generalization to sampling-based settings requires separate validation.

\textbf{Multi-sample baselines not in same pipeline.} The comparison to SE and SelfCheckGPT is cross-pipeline. The AUROC gaps are large, but within-pipeline comparison would be more definitive.

\textbf{TruthfulQA label protocol.} We use ROUGE-L $\geq 0.3$ for TruthfulQA binary labels (following \citet{fadeeva2024}). TruthfulQA is also evaluated with judge-model protocols; results may differ under alternative labeling schemes.

\subsection{Broader Impact}

This work advances zero-cost hallucination detection --- a practically important component of responsible LLM deployment. By providing a principled, label-free selection criterion for aggregation function, we reduce the risk of systematically underperforming detection for specific error classes. The mechanism-grounded taxonomy may generalize to other sequence-level tasks where different error types produce structurally distinct token probability distributions.
